%% ma: L12-B07-0017
%% lop: 12
%% chuong: 2
%% bai: 7
%% dang: 12.7.1
%% loai: TN4
%% muc: NB
%% dap_an: "D"
%% nguon: "CĐ2.2.pdf giáo viên gửi 27/10/2026 (M07, Lê Minh Khôi – Nguyễn Anh Khôi, Chương II Chủ đề 2 Tọa độ của vectơ trong không gian), Câu 14"
%% kien_thuc: "Bài 7 (hình chiếu của điểm lên mặt phẳng tọa độ)"
%% trang_thai: chinh-thuc
\begin{ex}
Trong không gian $Oxyz$, cho điểm $K(3;12;7)$. Gọi $K'$ là hình chiếu của $K$ lên mặt phẳng $(Oxy)$. Tọa độ của điểm $K'$ là
\choice{$K'(3;0;0)$}{$K'(0;12;0)$}{$K'(0;0;7)$}{\True $K'(3;12;0)$}
\loigiai{Hình chiếu lên $(Oxy)$ giữ hoành độ, tung độ và cho cao độ bằng $0$: $K'(3;12;0)$.}
\end{ex}
